\section{Related Work}
\label{sec:related-work}

\begin{table*}[!t]
\centering
\caption{Representative prior work and the gap this project addresses.}
\label{tab:prior-work}
\footnotesize
\setlength{\tabcolsep}{4pt}
\begin{tabular}{@{}>{\raggedright\arraybackslash\leavevmode}p{3.3cm}p{3.6cm}p{4.0cm}p{5.6cm}@{}}
\toprule
\textbf{Study} & \textbf{Focus} & \textbf{Method} & \textbf{Gap for this project} \\
\midrule
\citet{realinho2022predicting} & Dropout dataset (UCI) & Dataset description, ML baselines & Three-class accuracy; no capacity rule, explanation audit or users \\
\citet{qu2025tabicl}; \citet{hollmann2025tabpfn} & Tabular foundation models & In-context learning on tables & Benchmarks accuracy; not explanation quality or decision use \\
\citet{lou2013ebm}; \citet{caruana2015intelligible} & Interpretable models & Generalised additive models (EBM) & Glass-box alternative; predates, and is not compared with, foundation models \\
\citet{lundberg2017shap} & Local explanations & Shapley-value attributions & Approximations depend on budget and background data \\
\citet{tiukhova2024stability} & Learning analytics & Stability of XAI for student success & Stability of explanations, but no human evaluation \\
\citet{nauta2023evaluation} & XAI evaluation & Systematic review (Co-12 properties) & Calls for quantitative, multi-property evaluation \\
\citet{hase2020forward} & Human evaluation & Forward simulation & Not applied to dropout outreach \\
\citet{karimi2021recourse} & Recourse & Counterfactuals vs.\ interventions & Feasible changes need causal care \\
\midrule
\textbf{This project} & Dropout outreach at 20\% capacity & TabICL vs.\ 4 baselines, SHAP audit, 50-person study & Combines predictive, technical-explanation and human evidence \\
\bottomrule
\end{tabular}
\end{table*}

\subsection{Student Dropout Prediction}

The dataset used in this study~\citep{realinho2021uci} was released by \citet{realinho2022predicting},
who described 4,424 undergraduates from a Portuguese polytechnic through demographic,
socio-economic, application, first- and second-semester academic and macro-economic attributes,
labelled each student Dropout, Enrolled or Graduate, and reported three-class baselines with
standard classifiers. Such accuracy-oriented evaluations, however, use information from the whole
course of study and do not reflect the constraints of a support office. The present study departs
from that setting in three ways that mirror how such an office would use the scores: it fixes a
single prediction time, excludes later-semester information and evaluates performance at an
explicit outreach capacity.

\subsection{Tabular Foundation Models and Interpretable Baselines}

TabPFN~\citep{hollmann2025tabpfn} and \tabicl~\citep{qu2025tabicl} are transformers pretrained on
synthetic tables; they predict a new row by reading the training rows as context, without gradient
updates on the target data, and \tabicl extends this approach to larger tables through
column-then-row attention. Their reported strength on benchmarks makes them attractive for small
institutional datasets, yet their internal reasoning is not open to inspection. Strong conventional
baselines include logistic regression and gradient-boosted trees such as
XGBoost~\citep{chen2016xgboost}. Explainable Boosting Machines (EBMs), by contrast, are boosted
generalised additive models whose per-feature shape functions are themselves the
explanation~\citep{lou2013ebm, caruana2015intelligible}. \citet{rudin2019stop} argues that
interpretable models of this kind should be preferred in high-stakes settings when accuracy is
comparable, a claim that this study tests directly by comparing \tabicl with an EBM under the same
decision rule.

\subsection{Evaluating Explanations}

SHAP attributes a prediction to features using Shapley values~\citep{lundberg2017shap}. For
black-box models, these values are estimated by sampling and consequently depend on the sampling
budget, the random seed and the background data that defines the ``typical'' reference.
\citet{nauta2023evaluation} recommend measuring properties such as faithfulness and stability
quantitatively rather than relying on anecdotes, and \citet{tiukhova2024stability} show that
explanations of student-success models can be unstable across refits. Counterfactual
explanations~\citep{wachter2017counterfactual} describe which change would alter a prediction,
although \citet{karimi2021recourse} caution that a feasible change to a feature is not equivalent
to an effective real-world intervention. At the global level, model behaviour can be summarised with
accumulated local effects (ALE)~\citep{apley2020ale}, which remain reliable when features are
correlated. Taken together, these studies suggest that an explanation should be judged by several
measurable properties rather than by its plausibility alone.

\subsection{Human-Centred Evaluation}

\citet{doshivelez2018considerations} distinguish application-grounded, human-grounded and
functionally-grounded evaluation. Among human-grounded methods, forward simulation asks people to
predict a model's output from its input, with or without an explanation~\citep{hase2020forward}.
Such a test measures whether an explanation makes the model's behaviour predictable, not whether
people trust the model appropriately or make better decisions, and this distinction is maintained
throughout the present study.
